\newcommand{\DoNotLoadEpstopdf}{}
\documentclass[11pt,reqno]{amsart}
\usepackage[T1]{fontenc}
\usepackage{lmodern,amsmath,amssymb,amsthm,mathtools,mathrsfs,microtype}
\usepackage[margin=1.1in]{geometry}
\usepackage[hidelinks]{hyperref}
\hypersetup{pdftitle={Collision records, arithmetic, and raw local inversion in a triangular Lorentz gas},pdfauthor={Qian Qi}}
\setlength{\emergencystretch}{2em}
\allowdisplaybreaks[2]
\numberwithin{equation}{section}
\newtheorem{theorem}{Theorem}[section]
\newtheorem{lemma}[theorem]{Lemma}
\newtheorem{proposition}[theorem]{Proposition}
\newtheorem{corollary}[theorem]{Corollary}
\theoremstyle{definition}
\newtheorem{definition}[theorem]{Definition}
\newtheorem{inputstatement}[theorem]{Input statement}
\theoremstyle{remark}\newtheorem{remark}[theorem]{Remark}
\newcommand{\R}{\mathbb R}\newcommand{\Z}{\mathbb Z}
\newcommand{\one}{\mathbf1}\newcommand{\dd}{\,\mathrm d}
\newcommand{\esssup}{\operatorname*{ess\,sup}}
\newcommand{\diver}{\operatorname{div}}\newcommand{\TV}{\operatorname{TV}}
\title[Collision records and raw local inversion]{Collision records, arithmetic, and raw local inversion in a triangular Lorentz gas}
\author{Qian Qi}
\date{October 13, 2026. A2-DYN, revision 97}
\subjclass[2020]{37D50, 37A50, 60F05, 28A75}
\keywords{Lorentz gas, exact collision records, coarea occupation, maximal compression, original-source recovery, pointwise inversion}
\begin{document}
\raggedbottom
\begin{abstract}
We study the upper pointwise local limit problem for the exact record
of actual returns in a triangular Lorentz gas. Capping the physical
coarea ratio before time integration gives an occupation anchor using
both available directions and only returns to the identical selected
event. This recovers targets excluded by a whole-interval or worst-time
test. Independently, we bound the running compression of the complete
physical angular density at each target roof on prescribed same-event
intervals. The positive-cutoff inverse-gradient input pays this maximum
without counting words or visits. Together these estimates quantitatively
reduce the preceding three-way failure source to a subsource near an
original entrance with small gradient, or on a short original segment.
The artificial endpoint collar has no contribution at a uniform small
radius. Every recovered contribution has an explicit vanishing height
bound along the reconstruction schedule. The residual essential height
is not proved zero; the original pointwise theorem and all exact physical
and arithmetic data remain unchanged.
\end{abstract}
\maketitle
\input{journal-problem-v93}

\subsection{Quantitative recovery of failed targets}
The new results concern the original failed-test source, not only a
newly selected family of regular targets. With $Z_{r,D,\beta}^\varepsilon$
denoting the preceding disjoint interruption, compression and
occupation failures, Theorem~\ref{thm:v97-reduction} constructs an
original subsource $Y^\varepsilon\le Z^\varepsilon$ and proves
\[
 \mathcal H_M(Z^\varepsilon-Y^\varepsilon)
 \le Cq^{-1}\{\varepsilon^{1/16}+\varepsilon^{-1}\sqrt r\}+Cqr,
 \qquad q=\beta/D\le1/2.
\]
This is essential height at the target roof. It is not an integrated
translation or exceptional-set estimate. The residual support is
specified by the original lifetimes in \eqref{eq:v97-support}.
The running maximum responsible for the $Cqr$ term is estimated in
Theorem~\ref{thm:v97-running} before taking the collision limit.

All new theorems use the following explicitly retained inputs. The
comparison-domain anchor uses no source estimate outside the exact
event. The singular-volume estimate for running compression is applied
only while the trajectory remains in that event. This distinction
prevents label changes from entering a fixed-label spectral estimate.

\input{journal-inputs-v94}
\input{journal-comparison-v97}
\input{core/232_capped_coarea_occupation}
\input{core/233_running_compression_height}
\input{core/234_quantitative_failure_source_reduction}

\appendix
\input{appendices/v97_dependency_map}
\enlargethispage{3\baselineskip}
\input{journal-references-v93}
\end{document}
